\begin{abstract}
Android consumer applications (apps) provide everyday services such as news, social networking, messaging, and personalized content delivery, but these capabilities also introduce privacy and security exposure through permissions, exported components, storage access, embedded services, and network behavior. This paper presents an empirical assessment of privacy and security risk in Android consumer apps using ScytaleDroid, a database-backed Android analysis platform. The study evaluates a 15-app dataset using recent build-backed evidence bundles collected within a fixed 14-day cutoff window, preserving APK hashes, package versions, static run identifiers, dynamic run identifiers, PCAP availability, and quality-assurance provenance. Static analysis examines manifest declarations, permission posture, component exposure, detector findings, string and resource indicators, and MASVS-aligned risk areas across harvested base and split APK artifacts. Runtime analysis complements the static evidence with idle, quiescent foreground, and interactive captures collected from a non-root physical Android device. The results show that consumer apps expose measurable static risk surfaces while also exhibiting distinct runtime behavior patterns that are not visible from static evidence alone. The study further shows that app-update churn makes same-day installed-build completion impractical as an empirical reporting gate; recent cutoff-window evidence provides a more reproducible and auditable basis for analysis. In general, the paper demonstrates how database-backed evidence bundles improve repeatability, provenance, and interpretation in privacy and security studies.
\end{abstract}

\begin{IEEEkeywords}
Android security, mobile application privacy, static analysis, runtime analysis, dynamic analysis, APK analysis, MASVS, privacy risk, security risk, evidence provenance
\end{IEEEkeywords}
